\section{Introduction}\label{sec:introduction}

Consider the classical Li\'enard system
\begin{equation}\label{eq:original-system}
 \dot x=y-F(x),\qquad \dot y=-x,
\end{equation}
where $F$ is a real polynomial. A limit cycle is the image of a
nonconstant periodic solution that is isolated among periodic
orbits. We count each such image once, irrespective of its
multiplicity, stability, or hyperbolicity. The periodic orbits of
a center are not limit cycles under this convention.

\begin{theorem}\label{thm:main}
If $\deg F\le5$, the system \eqref{eq:original-system} has at most
two limit cycles in $\R^2$. Some members of this class have two
limit cycles. Thus its exact maximum is two.
\end{theorem}

Here the degree is that of the primitive $F$. The equivalent scalar
equation is
\[
 x''+F'(x)x'+x=0,
\]
so its damping polynomial has degree at most four and its restoring
force is exactly $x$. No parity or coefficient-sign assumption is
imposed on $F$.

The problem is a restricted instance of the second part of
Hilbert's sixteenth problem, concerning the number of limit cycles
of polynomial planar vector fields. Lins, de Melo, and Pugh proposed
the bound $\lfloor(n-1)/2\rfloor$ for
\eqref{eq:original-system} with $\deg F=n\ge1$ \cite{LMP}.
Rychkov proved the two-cycle bound for odd quintic primitives
\cite{Rychkov}. Dumortier, Panazzolo, and Roussarie constructed systems with
four cycles and primitive degree seven \cite{DPR}.
De Maesschalck and Dumortier obtained four cycles already in
degree six \cite{DMD}, and De Maesschalck and Huzak subsequently
constructed at least $n-2$ cycles in every degree $n\ge6$
\cite{DMH}. These counterexamples leave degree five undecided.
Li and Llibre established the unrestricted quartic bound of one
cycle \cite{LiLlibre}. Llibre and Zhang survey this development
and the classical perturbative lower bounds \cite{LlibreZhang}.

In degree five, the distinction between a global bound and a
bound in a limiting regime is essential. Li and Lu proved that
nondegenerate slow--fast cycles have cyclicity at most two
\cite{LiLu}. The Introduction of the August 2026 preprint by
Chen, Li, Zhang, and Zhang distinguishes that result from the
unrestricted polynomial problem and reports the latter as open
\cite{CLZZ}. Hern\'andez Rosales proposes a degree-five four-cycle
example for $x''+\mu f(x)x'+x=0$ in Section~9 of \cite{Rosales},
with $\mu=1/10$ and $f(x)=x^4-10x^2+5$. The proposed system has exactly two
limit cycles. Indeed, the change $x=\sqrt5\,u$ transforms its
scalar equation into
\[
 u''+\frac12(5u^4-10u^2+1)u'+u=0.
\]
This is Odani's Example~3 with parameters $\widehat\mu=1/2$ and
$v=10/3>\sqrt5$, for which exactly two periodic solutions are
proved \cite{Odani}. The discrepancy in the proposed amplitude
calculation is that periodicity requires
$\int_0^T f(x)\dot x^2\,dt=0$. Substituting
$x(t)\approx2X\cos t$ gives $2X^4-10X^2+5=0$, the same equation
used for the second proposed pair. Its roots are
$X^2=(5\pm\sqrt{15})/2$, not $1$ and $4$. Thus these two
leading-order calculations do not yield four distinct cycles.

Two-cycle examples are classical. Besides the perturbative
constructions associated with the conjecture \cite{LMP,LlibreZhang},
rigorous transversal-region methods have been applied to Rychkov's
family by Gasull, Giacomini, and Grau \cite{GGG}.
Section~\ref{sec:sharpness} supplies a direct lower-bound proof,
so the exact maximum follows within this article.

\paragraph{Method.}
The even coefficients prevent reduction to the odd quintic family,
and a small-amplitude calculation cannot by itself bound distant
cycles. Our comparison keeps the two half-orbits separate and
uses coordinates available at every transverse amplitude.
On each half-plane, $u=x^2/2$ transforms an orbit into an arch
satisfying $u_y=\phi(u)-y$, where
$\phi(u)=F(\sqrt{2u})$ or $F(-\sqrt{2u})$.
The even coefficients give the same quadratic part in both profiles,
whereas the odd coefficients change sign. This common part suggests
using quadratic profiles as a comparison family.
At a fixed base height $h$, let $y_<,y_>$ be its endpoints and set
$r=(y_>-y_<)/2$ and $M=(y_>+y_<)/2-\phi(h)$.
The first technical step is a global interpolation theorem:
every admissible pair $(M,M_r)$ has a unique fit by a profile
$\lambda(u-h)+\kappa(u-h)^2/2$. Both its range and its
nonvanishing Jacobian are proved in Theorem~\ref{thm:quadratic-fit};
the fit is available throughout $r>0$, $|M|<r$, and $|M_r|<1$.

The fitted slope and curvature evolve according to the two
transport equations of Lemma~\ref{lem:fit-transport}. Their
coefficients are derivatives of the quadratic midpoint function.
The central model estimate, Theorem~\ref{thm:model-inequality},
controls how one of these coefficients changes with width while
the model parameters are held fixed. Its proof combines endpoint
variation, a Riccati linearization, and a Schwarzian identity for
the endpoint correspondence. The solution-ratio identity and chain
rule are classical projective identities; Ovsienko and Tabachnikov
give a concise account \cite{OT}. Schwarzian derivatives have also
been used for cycle bounds through return maps, including the
discontinuous systems studied by Coll, Gasull, and Prohens
\cite{CGP}. Here the map pairs the two endpoints of a quadratic
arc, and the model estimate concerns that correspondence.

Transporting the fitted parameters along an arch then turns
conditions on $\phi'''$ and on $(\phi'-d)/u$, for a fixed real
constant $d$, into monotonicity and separation statements for the
fitted curvature and slope intercept. Propositions
\ref{prop:third-derivative} and~\ref{prop:intercept} formulate these
comparisons for general smooth profiles, independently of the
polynomial application.

For the original system, periodic orbits are zeros of the
difference between the two height-zero midpoint functions on
one common width interval. The quintic coefficient structure
provides a complete sign partition. In the only case allowing two
cycles, a positive integrating factor makes the derivative of
the matching function have the sign of a nondecreasing function.
This bounds isolated zeros even when zeros are multiple or occur
in intervals. The proof therefore counts degenerate cycles
directly; it does not remove them by a generic perturbation.

Section~\ref{sec:arcs} establishes the arc coordinates and the
global matching interpretation. Section~\ref{sec:fitting} proves the global quadratic fit and its
transport law. Section~\ref{sec:model} establishes the model
inequality, and Section~\ref{sec:transport} derives the shape
comparisons for general profiles. Section~\ref{sec:application} proves the upper bound in
Theorem~\ref{thm:main}, and Section~\ref{sec:sharpness} proves
sharpness.
